The fixed-field construction of Section~\ref{sec:limits} has
\[
 \begin{gathered}
 c_{\rm det}=a+4(h_K+2c_E+2t_K),\quad
 L_T(M)=16(d_K+5)M,\\
 f_x\widetilde u_x=d_x,\qquad
 c=v_2f_0(0)\le c_{\rm det}.
 \end{gathered}
\]
Choose $M>\max(c_{\rm det},B+2c_{\rm det},0)$, then
$b>(L_T(M)+2)(2^M-1)$ large enough for the address construction,
then its finite graph budget and actual theta stage $M_\theta$.
The residual comparison tests coefficients of $f$ from
$-L_T(M)$ through zero and only the constant coefficient of $d$.
It gives $x\ne0$ with
\[
       v_2(\widetilde u_x(0)-\widetilde u_0(0))\ge M-2c>B.
\]
The estimate uses both $v_{\mathscr O}(f_x)\le c$ and
$\widetilde u_x\in\Lambda[1/2]$; neither a bound on the whole
numerator nor a finite-Artin inverse of $t$ is asserted.
This contradicts (F.8) under the supposed base zero.
The Rankin factorization and $L(E^k,1)\ne0$ give the asserted
simple zero of $E$.

This completes the proof of the restricted converse,
Theorem~\ref{hyp:RC}, from the displayed classical inputs and
the arithmetic, coefficient and determinant constructions.
The conversion from that converse to pole witnesses, the integer
model, the natural compiler and finally the ordinary rational value
is the separate argument in the main text.

% End 09-normalization-endpoints.tex


